\documentclass[10pt,a4paper]{amsart}
\usepackage[margin=19mm]{geometry}
\usepackage{amsmath,amssymb,mathtools}
\usepackage[scr=rsfs,cal=euler]{mathalfa}
\usepackage{xfrac}
\usepackage[unicode,hidelinks,bookmarks=false]{hyperref}
\newcommand{\ie}{i.e.\ }
\newcommand{\cf}{cf.\ }
\newcommand{\resp}{respectively\ }
\newcommand{\Subsection}{\subsection}
\providecommand{\nobreakdash}{\nobreak\hbox{-}}
\setlength{\emergencystretch}{2em}
\multlinegap0pt
\usepackage{bm}
\usepackage{bbm}
\usepackage{dsfont}
\usepackage{xspace}
\usepackage{cancel}
\usepackage{mathtools}
\usepackage{cleveref}
\usepackage{amsthm}
\makeatletter
\@ifundefined{theorem}{\newtheorem{theorem}{Theorem}}{}
\makeatother
\title[A symbolic computational approach to the generalized gambler]{A symbolic computational approach to the generalized gambler's ruin problem in one and two dimensions}
\author{Lucy Martinez}
\date{Source: arXiv:2412.07667v1 (CC BY 4.0)}
\begin{document}
\maketitle

\noindent\emph{Source and licence.} A symbolic computational approach to the generalized gambler's ruin problem in one and two dimensions. Version arXiv:2412.07667v1 (CC BY 4.0), distributed under the Creative Commons Attribution 4.0 International licence, as recorded in the arXiv licence metadata for this exact version and shown on the abstract page https://arxiv.org/abs/2412.07667v1. Full rights notice: licenses/arxiv-2412.07667-CC-BY-4.0.txt.\par\smallskip

\noindent\textit{Editorial scope.} Counted unit: the Theorem whose source label is \texttt{thm:probability mirror} in the article (source0.tex lines 654--661, source label \texttt{thm:probability mirror}), delivered together with the article's own complete original proof of it (source0.tex lines 664--711, one \texttt{proof} environment of 48 lines). No mathematics is written, repaired or re-derived: statement and proof are the article's own text line for line. Required context retained uncounted: eq:inhomogeneous eq (lines 648--650). Every definition, display and label of those lines is retained unchanged. Omitted from this package and not redistributed: 1--647, 651--653, 662--663, 712--1107. Nothing inside a retained excerpt is omitted or altered: every retained body is delivered exactly as the article's lines are written, so no omission receipt of any kind accompanies this package. Citations: the retained excerpts carry no citation command at all, so no bibliography entry of the article is transcribed and no reference is redistributed. Reference adaptation: none was needed and none was made; every reference inside the delivered unit resolves inside it, so no unresolved reference or citation is produced. Printed slips kept literal: no printed word, symbol, hypothesis or number of the counted statement or of its proof was altered. Layout and class: the article's own class is \texttt{amsart} with packages including inputenc, xcolor, fullpage, hyperref, cleveref, enumitem, tikz, amsmath, amssymb, color, mathtools, float, pgf, amsrefs; the wrapper is the lane's pinned \texttt{amsart} editorial wrapper at 10pt on A4, which restates the theorem-like environments the retained text needs and adds the article's own macro definitions that the retained text needs (none), with extra wrapper packages (bm, bbm, dsfont, xspace, cancel, mathtools, cleveref). The delivered numbering is the unit's own. Assets: no retained span contains a figure environment, a graphics-inclusion command, a PDF-page inclusion, a table or a TikZ picture; no image file is copied into this package, the package declares no asset dependency and no image file is redistributed with it. Licence: version arXiv:2412.07667v1 of this article is distributed under the Creative Commons Attribution 4.0 International licence, as recorded in the arXiv licence metadata for this exact version and shown on the abstract page https://arxiv.org/abs/2412.07667v1. A full rights notice is delivered at licenses/arxiv-2412.07667-CC-BY-4.0.txt. Characters: every retained excerpt is pure ASCII, so the delivered engine, XeLaTeX with its default Latin Modern text font, reports no missing character.
\begin{equation}\label{eq:inhomogeneous eq}
    f(x)=\frac{p}{1+p}+\frac{1}{2}\left(\frac{1-p}{1+p}\right)f(x-1)+\frac{1}{2}\left(\frac{1-p}{1+p}\right)f(x+1),  \quad f(0)=0, f(N)=1.
\end{equation}

\begin{theorem}\label{thm:probability mirror}
Consider the generalization of the gambler's ruin problem when we add a mirror step. Then, the probability of ending at $N$ starting at $x$ is given by
\begin{equation*}
    f(x)=\frac{1}{2}\frac{\left(\frac{1-\sqrt{p}}{1+\sqrt{p}}\right)^N+1}{\left(\frac{1+\sqrt{p}}{1-\sqrt{p}}\right)^N-\left(\frac{1-\sqrt{p}}{1+\sqrt{p}}\right)^N}\left(\frac{1+\sqrt{p}}{1-\sqrt{p}}\right)^x
    +\frac{1}{2}\frac{\left(\frac{1+\sqrt{p}}{1-\sqrt{p}}\right)^N+1}{\left(\frac{1-\sqrt{p}}{1+\sqrt{p}}\right)^N-\left(\frac{1+\sqrt{p}}{1-\sqrt{p}}\right)^N}\left(\frac{1-\sqrt{p}}{1+\sqrt{p}}\right)^x+\frac{1}{2}
\end{equation*}
whenever we restrict the particle moves by either moving from $x$ to $x-1$ with probability $q_1$, or from $x$ to $x+1$ with probability $q_2$, or from $x$ to $N-x$ with probability $p$ where $q_1=q_2=\frac{1-p}{2}$.
\end{theorem}

\begin{proof}
    We can solve for \Cref{eq:inhomogeneous eq} because it is an inhomogeneous recurrence relation. Hence, we find a homogeneous and an inhomogeneous solution. 
    The general solution to the homogeneous equation $f(x)=\frac{1}{2}\left(\frac{1-p}{1+p}\right)f(x-1)+\frac{1}{2}\left(\frac{1-p}{1+p}\right)f(x+1)$ is
    
    \begin{equation*}
        f(x)=A\left(\frac{1+p+2\sqrt{p}}{1-p}\right)^x+B\left(\frac{1+p-2\sqrt{p}}{1-p}\right)^x
    \end{equation*}
    
    for some numbers $A$ and $B$. 
    Next, we find the particular solution to the inhomogeneous relation by setting $f^*(x)=C$ for some constant $C$. Then,
    
    \begin{equation*}
        f^*(x)=\frac{p}{1+p}+\frac{1}{2}\left(\frac{1-p}{1+p}\right)f^*(x-1)+\frac{1}{2}\left(\frac{1-p}{1+p}\right)f^*(x+1)
    \end{equation*}
    
    becomes
    
    \begin{equation*}
        C=\frac{p}{1+p}+\frac{1}{2}\left(\frac{1-p}{1+p}\right)C+\frac{1}{2}\left(\frac{1-p}{1+p}\right)C
    \end{equation*}
    
    which has solution $C=\frac{1}{2}$. 
    
    Therefore, $f^*(x)=\frac{1}{2}$ is the particular solution, and the general inhomogeneous solution is
    
    \begin{equation}\label{eq:general inhomogeneous sln}
         f(x)=A\left(\frac{1+p+2\sqrt{p}}{1-p}\right)^x+B\left(\frac{1+p-2\sqrt{p}}{1-p}\right)^x+\frac{1}{2}.
    \end{equation}
    
    Using Maple, we find $A$ and $B$ by using the boundary conditions to get a system of two linear equations. Namely,
    
    \begin{align*}
        0&=A+B+\frac{1}{2} \intertext{and}
        \frac{1}{2}&=A\left(\frac{1+p+2\sqrt{p}}{1-p}\right)^N+B\left(\frac{1+p-2\sqrt{p}}{1-p}\right)^N.
    \end{align*}
    
   Solving for the above linear equations, we get
   \begin{align}
       A&=\frac{1}{2}\frac{\left(\frac{1+p-2\sqrt{p}}{1-p}\right)^N+1}{\left(\frac{1+p+2\sqrt{p}}{1-p}\right)^N-\left(\frac{1+p-2\sqrt{p}}{1-p}\right)^N} \label{eq: A eq}\intertext{and}
       B&=\frac{1}{2}\frac{\left(\frac{1+p+2\sqrt{p}}{1-p}\right)^N+1}{\left(\frac{1+p-2\sqrt{p}}{1-p}\right)^N-\left(\frac{1+p+2\sqrt{p}}{1-p}\right)^N}.\label{eq: B eq}
   \end{align}
   
   Rewriting Equations \ref{eq:general inhomogeneous sln}, \ref{eq: A eq} and \ref{eq: B eq} gives
   \begin{align*}
       f(x)&=A\left(\frac{1+\sqrt{p}}{1-\sqrt{p}}\right)^x+B\left(\frac{1-\sqrt{p}}{1+\sqrt{p}}\right)^x+\frac{1}{2} 
   \end{align*}
as desired.
\end{proof}

\end{document}
